\documentclass[11pt]{article}

\usepackage{amsmath,amssymb}
\usepackage{xurl}
\usepackage[hidelinks]{hyperref}
\setlength{\emergencystretch}{2em}

% Shared commands for notes in docs/paper-gaps/.
% Notation follows the LDT paper (references/ldt-paper/).

\newcommand{\Lean}{\textsc{Lean}}
\newcommand{\leanid}[1]{\textsf{#1}}
\newcommand{\ghissue}[1]{\href{https://github.com/LionSR/MIPStarRE/issues/#1}{GitHub issue \##1}}
\newcommand{\ghpr}[1]{\href{https://github.com/LionSR/MIPStarRE/pull/#1}{GitHub PR \##1}}

% --- Paper-compatible notation ---
\newcommand{\F}{\mathbb{F}}
\newcommand{\N}{\mathbb{N}}
\newcommand{\C}{\mathbb{C}}
\newcommand{\tr}{\operatorname{tr}}
\newcommand{\ot}{\otimes}
\newcommand{\E}{\mathop{\bf E\/}}
\newcommand{\bx}{\mathbf{x}}
\newcommand{\by}{\mathbf{y}}
\newcommand{\bu}{\mathbf{u}}
\newcommand{\bv}{\mathbf{v}}

% Approximation relation (paper uses \approx_\eps and \simeq_\zeta)
\newcommand{\approxeps}{\approx_{\varepsilon}}
\newcommand{\simgeq}{\simeq_{\zeta}}

% Polynomial spaces (paper uses \polyfunc, \polysub, \polymeas)
\newcommand{\polyfunc}[3]{\operatorname{Poly}(#1,#2,#3)}
\newcommand{\polysub}[3]{\operatorname{PolySub}(#1,#2,#3)}
\newcommand{\polymeas}[3]{\operatorname{PolyMeas}(#1,#2,#3)}


\title{Issue 196: The SDP Primal Constraint Is a Submeasurement Constraint}
\author{}
\date{2026-04-30}

\newcommand{\Tr}{\operatorname{Tr}}
\newcommand{\cP}{\mathcal P}

\begin{document}
\maketitle

\section{At a glance}

\begin{itemize}
  \item \textbf{Difficulty: easy.} This issue is about the feasible set of one
  displayed SDP, not about proving SDP strong duality. The strong-duality and
  complementary-slackness producer is documented separately.%
  \footnote{The full strong-duality producer is discussed in
  \path{docs/paper-gaps/issue-1230-self-improvement-sdp-usage.tex}.}
  \item \textbf{Estimated weight:} already resolved; replaying the repair is about
  \textbf{2--4 interface changes} and roughly \textbf{100--250 lines} of
  project-local definitions and adapters.
  \item \textbf{Mathlib/project split:} about \textbf{20\% Mathlib} and
  \textbf{80\% project-local}. Mathlib supplies finite sums and matrix order;
  the project supplies the submeasurement, measurement, SDP-wrapper, and
  completion interfaces.
  \item \textbf{Key Mathlib inputs:} finite sums over the polynomial set, matrix
  positive-semidefinite/order facts, and scalar multiplication of positive
  operators.
\end{itemize}

\section{Key theorem forms}

The resolved formal interface should isolate four statements.

\begin{enumerate}
  \item \textbf{Submeasurement primal form.} The primal variable is a family
  $T=\{T_g\}$ with $T_g\ge0$ and $\sum_gT_g\le I$.
  \item \textbf{Strict feasible witness.} For $M=|\cP|$, the family
  $T_g=(2M)^{-1}I$ satisfies $\sum_gT_g=\frac12 I$ and is strictly feasible.
  \item \textbf{Saturation as output.} The identity $\sum_gT_g=I$ belongs to the
  selected optimal witness, together with $T_gZ=T_gA_g$; it is not a feasibility
  condition.
  \item \textbf{Completion only downstream.} If a later wrapper needs a full
  measurement, complete the chosen submeasurement explicitly, and do not identify
  that completion with the paper's strict Slater witness.
\end{enumerate}


\section{The primal feasible set in the paper}

This note concerns the semidefinite program used in the self-improvement
argument of \cite[Subsection~\emph{A semidefinite program}]{Ji2020LowIndividualDegree}.
\footnote{The mismatch was recorded in
\href{https://github.com/LionSR/MIPStarRE/issues/196}{GitHub issue \#196} and
repaired in
\href{https://github.com/LionSR/MIPStarRE/pull/251}{GitHub pull request \#251}.
The documentation issue for this note is
\href{https://github.com/LionSR/MIPStarRE/issues/935}{GitHub issue \#935}. In
the TeX source consulted here, the primal SDP is in
\path{references/ldt-paper/self_improvement.tex}, lines 64--74; the SDP lemma is
in lines 82--88; and the Slater witness is in lines 168--176.}
Let $\cP$ be the finite set of degree-bounded polynomials, and set
\[
  A_g = \mathbb E_{u\in \mathbb F_q^m} A^u_{g(u)}
  \qquad (g\in\cP).
\]
The primal SDP in the paper is
\begin{align}
  \sup_{\{T_g\}_{g\in\cP}}
    &\quad \sum_{g\in\cP} \Tr(T_g A_g) \tag{\texttt{eq:primal-objective}}\\
  \text{subject to}
    &\quad T_g\ge 0 \qquad(g\in\cP), \notag\\
    &\quad \sum_{g\in\cP} T_g\le I. \notag
\end{align}
Thus a feasible primal variable is a submeasurement: a family of positive
semidefinite effects whose total mass is bounded above by the identity. It is
not, at this point, required to be a measurement.

The same lemma later asserts the existence of an optimal primal-dual pair
$(\{T_g\},Z)$ for which
\[
  \sum_{g\in\cP} T_g=I,
  \qquad
  T_g Z=T_g A_g \quad(g\in\cP).
  \tag{\texttt{lem:sdp}}
  \label{lem:sdp}
\]
This equality is a conclusion about a chosen optimal solution. It is not the
feasibility constraint of the primal SDP.

\section{Why equality is a stronger formulation}

Replacing $\sum_g T_g\le I$ by $\sum_g T_g=I$ changes the feasible set. The
change is proper. If $M=|\cP|$, the paper's strict primal witness is
\[
  T_g=\frac{1}{2M}I \qquad(g\in\cP).
\]
Its total mass is
\[
  \sum_{g\in\cP} T_g=\frac12 I,
\]
so it is feasible for the paper's SDP and satisfies the inequality strictly, but
it is not feasible for the equality-constrained SDP.

This point matters for the proof. The family above is the strict-feasibility
witness used to invoke Slater's condition. It is deliberately not saturated:
the missing mass is represented in the canonical SDP by the additional slack
block $X_{M+1,M+1}$. Complementary slackness then forces this slack block to
vanish at an optimal solution, which yields $\sum_g T_g=I$ for the selected
optimum. If equality is imposed at the start, this derivation is replaced by an
assumption and the strict feasible point used in the paper is excluded.

Thus the equality formulation is not a harmless restatement of the displayed
primal program. It is a smaller SDP. It may still contain saturated optima, but
that saturation is exactly part of what the paper proves from the submeasurement
program.

\section{The current formal encoding}

The current development has a public paper-local type for submeasurements: a
family $T=\{T_g\}$ with $T_g\ge 0$, total
$\sum_g T_g$, and proof that this total is at most $I$. A measurement is the
special case where the total is exactly $I$.
\footnote{These are the declarations
\leanid{MIPStarRE.LDT.SubMeas} and \leanid{MIPStarRE.LDT.Measurement} in
\path{MIPStarRE/LDT/Basic/SubMeasurementCore.lean}, lines 13--28.}
The self-improvement SDP definitions now use the submeasurement type for the
primal variable. In particular, the primal contribution, primal objective, and
complementary-slackness equation all take $T$ to be a submeasurement, not a
measurement.
\footnote{See
\leanid{MIPStarRE.LDT.SelfImprovement.sdpPrimalContributionOperator},
\leanid{MIPStarRE.LDT.SelfImprovement.sdpPrimalObjectiveOperator},
\leanid{MIPStarRE.LDT.SelfImprovement.sdpPrimalObjective}, and
\leanid{MIPStarRE.LDT.SelfImprovement.sdpComplementarySlacknessEquation} in
\path{MIPStarRE/LDT/SelfImprovement/Defs.lean}, lines 139--184. The same file
also defines the sandwiched family \path{sandwichedPolynomialSubMeasAt} and the
averaged family \path{averagedSandwichedPolynomialSubMeas} with
$T$ a submeasurement, in lines 198--333.}
The proof that the sandwiched family is a submeasurement uses the inequality
$\sum_g T_g\le I$: after grouping the summands by the value of $h(u)$, the
filtered sum is bounded by the total mass of $T$, and hence by $I$.

The strict primal witness from the paper is represented before completion as a
submeasurement. The declaration \path{sdpStrictPrimalSubMeas} has outcomes
$T_g=(2|\cP|)^{-1}I$, and the theorem
\path{sdpStrictPrimalSubMeas_total} proves that its total is $\frac12 I$.
\footnote{These are in
\path{MIPStarRE/LDT/SelfImprovement/Defs.lean}, lines 37--97.}
This is the formal counterpart of the paper's Slater witness. Earlier Lean
versions also exposed a declaration \leanid{sdpPrimalWitness}, which completed
this submeasurement at the zero polynomial for a reduced theorem named
\leanid{sdp}.  That completion wrapper and the reduced theorem it served were
retired in \ghpr{2430}.  The active SDP route keeps the strict witness as a
submeasurement and obtains a saturated measurement only from the optimal-pair
conclusion.  Thus completion is not part of the primal feasible set of the SDP.

The informal blueprint already agrees with the paper on this point: its SDP
lemma states the primal with $\sum_g T_g\le I$, and only afterwards states the
existence of an optimal pair whose total is $I$. The resolved mismatch was
therefore confined to the formal encoding, not a three-way disagreement with the
blueprint.
\footnote{See \path{blueprint/src/chapter/ch07_self_improvement.tex}, lines
281--304.}

At the statement layer, the active source-facing route exposes the primal
witness as a measurement inside
\leanid{SdpStatementWithSlackness.witness}. The associated witness record is
indexed by the underlying submeasurement and still contains a field
\path{primalTotalOperator} asserting that this total is $I$; this is the
paper's saturated-optimum conclusion, not a definition of primal feasibility.
\footnote{See
\leanid{MIPStarRE.LDT.SelfImprovement.SdpOptimalPairWithSlackness} and
\leanid{MIPStarRE.LDT.SelfImprovement.SdpStatementWithSlackness} in
\path{MIPStarRE/LDT/SelfImprovement/Theorems/Statements.lean}, lines 61--134.
The source-facing theorem is
\leanid{MIPStarRE.LDT.SelfImprovement.sdp_statement_with_slackness} in
\path{MIPStarRE/LDT/SelfImprovement/Theorems/Results/HelperCompleteness/Bracketed.lean},
lines 493--510, using the canonical block-SDP transport from
\path{MIPStarRE/LDT/SelfImprovement/Theorems/Results/SdpMatrixBridge.lean}.}
The matrix-level realization follows the same convention: primal variables are
matrix submeasurements, while the slackness statement stores the completed
matrix measurement obtained from the identity-total conclusion.
\footnote{The matrix-level primal-side declarations
\path{matrixSdpPrimalContributionOperator}, \path{matrixSdpPrimalObjective},
\path{matrixSdpComplementarySlacknessDefect}, and \path{MatrixSdpOptimalWitness}
are in \path{MIPStarRE/LDT/SelfImprovement/MatrixRealization.lean}, lines
42--96; the witness structure itself is in lines 81--96.}

\section{Conclusion}

The verdict is that
\href{https://github.com/LionSR/MIPStarRE/issues/196}{issue \#196} recorded a
genuine mismatch. The paper's primal SDP is over submeasurements satisfying
$\sum_g T_g\le I$. The earlier equality-based formal encoding treated the
primal variable as a measurement and thereby excluded the paper's strict-feasible
Slater witness $T_g=(2|\cP|)^{-1}I$.

The mismatch is now resolved at the SDP-interface level. The formal objective,
complementary-slackness relation, strict primal witness, and sandwiched-family
construction all expose the feasible primal variable as a submeasurement.
Saturation is recorded only for the selected optimal witness, and the
slackness-carrying SDP statement stores that saturated witness as a measurement.

The former reduced interface has been retired: \ghpr{2430} removed
\leanid{sdpPrimalWitness}, the reduced \leanid{SdpStatement} structure, the theorem
named \leanid{sdp}, and the wrapper \leanid{selfImprovementHelperConstruction}. The
current \Lean{} route for Lemma~\ref{lem:sdp} is
\leanid{sdp_statement_with_slackness}, whose output is
\leanid{SdpStatementWithSlackness}.  The canonical block-SDP strong-duality route
is documented in
\path{docs/paper-gaps/issue-1230-self-improvement-sdp-usage.tex}. Future
refactors should preserve the submeasurement feasible set until the point where
the paper derives the saturated optimal witness $\sum_g T_g=I$.

\bibliographystyle{alpha}
\bibliography{references}

\end{document}
